\subsubsection{Tập dữ liệu}

Nghiên cứu này sử dụng tập dữ liệu UA-DETRAC (University at Albany DEtection
and TRACking) làm cơ sở chuẩn đối sánh. Quy mô của tập dữ liệu bao gồm 100
chuỗi video được chọn từ hơn 10 giờ từ camera ghi lại các bối cảnh giao thông
thực tế ở 24 vị trí khác nhau với đa dạng mẫu và điều kiện giao thông bao gồm
cao tốc đô thị, nút giao thông, ngã ba chữ T~\cite{wen2020uadetrac}. Các video
này được ghi lại với 25 khung hình/giây (FPS) với định dạng JPEG và độ phân
giải $960 \times 540$~\cite{wen2020uadetrac}.

Hơn 140,000 khung hình được chú thích (annotated) thủ công với hơn 8,250 phương
tiện và tổng cộng 1.21 triệu hộp bao được gán nhãn. Mỗi mẫu dữ liệu chứa thông
tin tọa độ của các hộp bao nhằm định vị chính xác các phương tiện. Bên cạnh đó,
mỗi đối tượng trong mẫu được mô tả chi tiết thông qua các thuộc tính bao gồm:
phân lớp phương tiện (ô tô, xe buýt, xe tải nhỏ và các loại khác), bối cảnh
chiếu sáng (nhiều mây, ban đêm, trời nắng, trời mưa), phạm vi kích thước (nhỏ,
trung bình, lớn), tỷ lệ che khuất (không che khuất, che khuất một phần, che
khuất nhiều) và tỷ lệ nằm ngoài khung hình~\cite{wen2020uadetrac}. Nhằm phục vụ
cho quá trình đánh giá mô hình, tập dữ liệu được phân hoạch thành hai tập con
độc lập: tập huấn luyện (train) chứa 60 chuỗi video với xấp xỉ 84.000 khung
hình, và tập kiểm tra (test) bao gồm 40 chuỗi video với 56.000 khung
hình~\cite{wen2020uadetrac}.
